\section{A reduction for categories of embeddings}\label{sec:reduction}

Both examples of this paper are categories of finite substructures of a
homogeneous structure. This section isolates what is used of such a category.

\subsection{The setting}

Let $M$ be a non-empty set and $G$ a group of permutations of $M$. For
$X\subseteq M$ let $G_X$ be the pointwise stabilizer of $X$; so
$G_\varnothing=G$. A subgroup of $G$ is \emph{open} if it contains $G_X$ for
some finite $X\subseteq M$. Let $\K$ be a set of finite subsets of $M$. We
consider the following conditions.
\begin{enumerate}
\item[\textup{(K1)}] $\varnothing\in\K$; $gB\in\K$ for $B\in\K$ and $g\in G$; every
 finite subset of $M$ is contained in a member of $\K$.
\item[\textup{(K2)}] $\Fix_M(G_B)=B$ for every $B\in\K$.
\item[\textup{(K3)}] Every open subgroup of $G$ is of the form $G_F$ with $F\in\K$.
\end{enumerate}
Let $\mathcal A=\mathcal A(M,G,\K)$ be the category whose objects are the
non-empty members of $\K$ and whose arrows $B\to C$ are the maps
$g|_B\colon B\to C$ with $g\in G$ and $gB\subseteq C$, composed as maps.
Inclusions between members of $\K$ are arrows. By (K1) any two objects are
contained in a third, so $\mathcal A$ is non-empty and connected.

For an object $P$ of a category $\mathcal C$, the objects of
$P\backslash\mathcal C$ are the pairs $(X,x\colon P\to X)$, and an arrow
$(X,x)\to(Y,y)$ is an arrow $f\colon X\to Y$ with $fx=y$. A \emph{cosieve} of a
category is a full subcategory that contains the target of every arrow whose
source it contains.
\begin{enumerate}
\item[\textup{(SC)}] For every object $P$ of $\mathcal C$, every non-empty cosieve
 of $P\backslash\mathcal C$ is weakly contractible.
\item[\textup{(K4)}] The category $\mathcal A(M,G,\K)$ satisfies \textup{(SC)}.
\end{enumerate}
Property (SC) is invariant under equivalence of categories: an equivalence
$\Phi\colon\mathcal C\to\mathcal C'$ induces equivalences
$P\backslash\mathcal C\to\Phi P\backslash\mathcal C'$, and the objects
isomorphic to images of objects of a cosieve $\mathcal U$ form a cosieve
equivalent to $\mathcal U$.

\begin{lemma}\label{lem:squares}
In a category with \textup{(SC)} every pair of arrows $f\colon a\to b$,
$f'\colon a\to b'$ can be completed to a commutative square $uf=u'f'$.
\end{lemma}

\begin{proof}
In $a\backslash\mathcal C$ let $\mathcal U$ be the cosieve of the objects
$(c,h)$ such that $h$ factors through $f$, and $\mathcal U'$ the one defined by
$f'$. Their union is a non-empty cosieve, hence connected. A path in it from
$(b,f)$ to $(b',f')$ contains an arrow between an object of $\mathcal U$ and an
object of $\mathcal U'$, and its target lies in $\mathcal U\cap\mathcal U'$.
\end{proof}

\subsection{The components of the categories of elements}

\begin{proposition}\label{prop:reduction}
Assume \textup{(K1)--(K4)}, and let $L\colon\mathcal A\to\Set$ be a functor. For
every connected component $D$ of $\El(L)$, either $D$ is weakly contractible,
or the projection $D\to\mathcal A$ is a weak equivalence.
\end{proposition}

\begin{proof}
A functor that induces a weak equivalence of classifying spaces induces a
bijection between the sets of connected components and weak equivalences
between corresponding components. The two functors constructed below commute
with the projections to $\mathcal A$.

\emph{Step 1: injective transition maps.} For $l,l'\in L(a)$ write $l\sim l'$
if $L(f)l=L(f)l'$ for some arrow $f\colon a\to b$. This is an equivalence
relation: if $f$ and $f'$ witness $l\sim l'$ and $l'\sim l''$, and $uf=u'f'$ as
in Lemma~\ref{lem:squares}, then $L(uf)$ takes the same value on $l$, $l'$ and
$l''$. It is preserved by $L(v)$ for every arrow $v\colon a\to a'$: if
$L(f)l=L(f)l'$ and $uf=u'v$, then $L(u')$ identifies $L(v)l$ and $L(v)l'$.
So the quotients form a functor $L^{\mathrm{inj}}$, with a natural surjection
$L\to L^{\mathrm{inj}}$. Its transition maps are injective: an arrow that
witnesses $L(v)l\sim L(v)l'$, composed with $v$, witnesses $l\sim l'$.

We claim that $q\colon\El(L)\to\El(L^{\mathrm{inj}})$ is a weak equivalence. Fix
an object $(a,[l])$ of the target and let $\mathcal D=(a,[l])\backslash q$; its
objects are the triples $(x,k,f)$ with $f\colon a\to x$, $k\in L(x)$ and
$k\sim L(f)l$. Let $J\colon a\backslash\mathcal A\to\mathcal D$ send
$f\colon a\to x$ to $(x,L(f)l,f)$. For an object $d=(x,k,f)$ of $\mathcal D$,
the category $d\backslash J$ is the full subcategory of $x\backslash\mathcal A$
on the arrows $g\colon x\to y$ with $L(g)k=L(gf)l$. It is a cosieve, and it is
non-empty because $k\sim L(f)l$; by (K4) it is weakly contractible. By
Quillen's Theorem~A \cite[\S 1]{Quillen}, $J$ is a weak equivalence. As
$a\backslash\mathcal A$ has an initial object, $\mathcal D$ is weakly
contractible, and Theorem~A applies to $q$.

\emph{Step 2: saturation.} By Step 1 we may assume that all maps $L(f)$ are
injective. The non-empty members of $\K$ form a directed set under inclusion,
by (K1). Let
\[
 S=\colim_{B}L(B)
\]
be the colimit over the inclusions. The canonical maps $i_B\colon L(B)\to S$
are injective. The group $G$ acts on $S$ by
$g\cdot i_B(l)=i_{gB}\bigl(L(g|_B)l\bigr)$, where $g|_B\colon B\to gB$; this is
well defined because $g|_{B'}$ extends $g|_B$ when $B\subseteq B'$. The element
$i_B(l)$ is fixed by $G_B$. So $S$ is a $G$-set all of whose point stabilizers
are open.

For any such $G$-set $S$ define a functor $T_S\colon\mathcal A\to\Set$ by
$T_S(B)=S^{G_B}$ and $T_S(g|_B)(s)=g\cdot s$ for an arrow $g|_B\colon B\to C$.
This does not depend on $g$, since two choices differ by an element of $G_B$;
and $g\cdot s$ is fixed by $G_C$ because $g^{-1}G_Cg\le G_B$. In our case the
maps $i_B$ form a natural monomorphism $i\colon L\to T_S$. For an object
$(a,s)$ of $\El(T_S)$, the category $(a,s)\backslash\El(i)$ is the full
subcategory of $a\backslash\mathcal A$ on the arrows $f\colon a\to b$ with
$T_S(f)s\in i_b(L(b))$. It is a cosieve. It is non-empty: if $s=i_c(l)$,
choose $b\in\K$ containing $a\cup c$; for the inclusion $f\colon a\to b$ we
have $T_S(f)s=s\in i_b(L(b))$. By (K4) and Theorem~A,
$\El(i)\colon\El(L)\to\El(T_S)$ is a weak equivalence.

\emph{Step 3: orbits.} The functor $S\mapsto T_S$ takes disjoint unions of
$G$-sets to disjoint unions of functors, so $\El(T_S)$ is the disjoint union of
the categories $\El(T_O)$, $O$ running through the orbits of $S$. By (K3)
every orbit is isomorphic to $G/G_F$ for some $F\in\K$.

If $F\ne\varnothing$, then $T_{G/G_F}\cong\mathcal A(F,-)$. Indeed, a coset
$gG_F$ is fixed by $G_B$ if and only if $G_B\le gG_Fg^{-1}=G_{gF}$, that is, by
(K2), if and only if $gF\subseteq B$. The map $gG_F\mapsto g|_F$ is then a
bijection from $(G/G_F)^{G_B}$ onto $\mathcal A(F,B)$, natural in $B$. The
category of elements of $\mathcal A(F,-)$ is $F\backslash\mathcal A$, which has
an initial object.

If $F=\varnothing$, the orbit is a point, $T_O$ is the constant functor with
value a point, and $\El(T_O)\to\mathcal A$ is an isomorphism.

So every component of $\El(T_S)$ is contractible or is mapped isomorphically
to $\mathcal A$ by the projection, and the statement follows from Steps 1
and 2.
\end{proof}

\subsection{The bridge set}

The following construction is used twice, to verify \textup{(K3)}. It involves
only finite linearly ordered sets.

\begin{lemma}\label{lem:bridgeset}
Let $D$ and $E$ be finite linearly ordered sets, $B,C\subseteq D$,
$F=B\cap C$, and let $d_0,d_1\colon D\to E$ be increasing injections that agree
on $F$. For every sufficiently large integer $N\ge 2$, with $s=2N$, there are a
finite linearly ordered set $Z$, an increasing injection $e\colon E\to Z$ and
increasing injections $a_0,\dots,a_s\colon D\to Z$ with the following
properties. Put $Z_0=e(E)$ and $Z_i=a_i(D)$ for $0<i<s$; these are the
\emph{pieces}, and their indices are read modulo $s$.
\begin{enumerate}
\item[\textup{(i)}] $Z=\bigcup_iZ_i$; $a_0=ed_0$ and $a_s=ed_1$;
 $a_{2j}|_B=a_{2j+1}|_B$ and $a_{2j+1}|_C=a_{2j+2}|_C$ for $0\le j<N$.
\item[\textup{(ii)}] For $x\in F$ the point $a_i(x)$ does not depend on $i$; these
 \emph{$F$-points} lie in every piece. Every other point of $Z$ lies in at
 most three pieces. If it lies in two pieces, these are consecutive, or they
 are $Z_1$ and $Z_{s-1}$, and then the point lies in $Z_0$ as well.
\item[\textup{(iii)}] For $0\le i<s$ we have $Z_i\cap Z_{i+1}=a_i(B)=a_{i+1}(B)$ if
 $i$ is even and $Z_i\cap Z_{i+1}=a_i(C)=a_{i+1}(C)$ if $i$ is odd.
\item[\textup{(iv)}] For every index $r$ there are a finite set $L$, a
 surjection $q\colon L\to Z$ and subsets $T_0,\dots,T_{s-1}$ covering $L$ such
 that $q$ maps $T_k$ bijectively onto $Z_{r+k}$, with these properties: every
 point of $L$ lies in $T_k$ for a set of consecutive values of $k$, with at
 most three elements unless the point lies over an $F$-point;
 $q(T_{k-1}\cap T_k)=Z_{r+k-1}\cap Z_{r+k}$ for $0<k<s$; every $F$-point has
 one preimage under $q$; and $q$ is injective except that at most
 $\max(|B|,|C|)$ points of $Z$ have two preimages.
\end{enumerate}
\end{lemma}

\begin{proof}
Let $\widehat Z$ be the disjoint union of $E$ and of copies $D_1,\dots,D_{s-1}$
of $D$; write $x_i$ for the copy of $x\in D$ in $D_i$. We call $E$ the
\emph{piece} $0$ and $D_i$ the piece $i$. Consider the following \emph{links}
between points of consecutive pieces: $d_0(x)$ and $x_1$, for $x\in B$; $x_i$
and $x_{i+1}$, for $0<i<s-1$ and $x\in B$ if $i$ is even, $x\in C$ if $i$ is
odd; $x_{s-1}$ and $d_1(x)$, for $x\in C$. Let $Z$ be the quotient of
$\widehat Z$ by the equivalence relation generated by the links, $e$ the
restriction of the quotient map to $E$, $a_i(x)$ the class of $x_i$ for
$0<i<s$, and $a_0=ed_0$, $a_s=ed_1$. Then (i) holds.

A point has at most one link to the next piece and at most one to the
preceding piece, because $d_0$ and $d_1$ are injective. For $x\in F$ the
points $d_0(x)=d_1(x),x_1,\dots,x_{s-1}$ form a class, with one point in each
piece. For $x\in B\setminus C$ the links at copies of $x$ join $x_i$ to
$x_{i+1}$ for $i$ even, $0<i<s-1$, and $d_0(x)$ to $x_1$; for
$x\in C\setminus B$ they join $x_i$ to $x_{i+1}$ for $i$ odd, $0<i<s-1$, and
$x_{s-1}$ to $d_1(x)$; for $x\notin B\cup C$ there are none. A point of $E$ that
is not of the form $d_0(x)$, $x\in F$, has at most two links, to some $x_1$
with $x\in B\setminus C$ and to some $y_{s-1}$ with $y\in C\setminus B$. Hence a
class that does not come from $F$ is a single point, or a pair
$\{x_i,x_{i+1}\}$, or is contained in a set $\{y_{s-1},d,x_1\}$ with $d\in E$.
So every class has at most one point in each piece, which shows that $e$ and
the $a_i$ are injective, and (ii) and (iii) follow by inspection; here
$N\ge 2$ is used.

For (iv), let $L$ be the quotient of $\widehat Z$ by the relation generated
by the links other than those between the pieces $r-1$ and $r$, let $T_k$ be
the image of the piece $r+k$, and $q$ the quotient map. A class of $Z$ is the
union of at most two classes of $L$, and of two only if it contains a link
between the pieces $r-1$ and $r$ and does not come from $F$; there are at
most $\max(|B|,|C|)$ such links. The remaining assertions follow from the
description of the classes.

\emph{The order.} Choose an increasing injection $\varepsilon\colon E\to\R$
and put $p_t=(1-t)\,\varepsilon d_0+t\,\varepsilon d_1\colon D\to\R$ for
$0\le t\le 1$. Each $p_t$ is strictly increasing, and $p_t(x)$ does not
depend on $t$ for $x\in F$. Let $\Delta>0$ be a lower bound for
$p_t(x')-p_t(x)$ over all $t$ and all $x<x'$ in $D$, and let
$M_0=\max_x|\varepsilon d_1(x)-\varepsilon d_0(x)|$; assume $2M_0<N\Delta$.
Define $b_i\colon D\to\R$ for $0\le i\le s$ by
\[
 b_{2j}=p_{j/N},\qquad
 b_{2j+1}(x)=\begin{cases}p_{j/N}(x),&x\in B,\\ p_{(j+1)/N}(x),&x\notin B.\end{cases}
\]
Each value of $b_{2j+1}$ differs from the corresponding value of $p_{(j+1)/N}$
by at most $M_0/N$, so $b_{2j+1}(x')-b_{2j+1}(x)\ge\Delta-2M_0/N>0$ for
$x<x'$: all $b_i$ are strictly increasing. Moreover $b_0=\varepsilon d_0$,
$b_s=\varepsilon d_1$, $b_{2j}$ and $b_{2j+1}$ agree on $B$, and $b_{2j+1}$
and $b_{2j+2}$ agree on $C$ (on $C\setminus B$ by definition, on $F$ because
$p_t$ is constant there). So the function on $\widehat Z$ that is
$\varepsilon$ on $E$ and $b_i$ on $D_i$ takes the same value at the two ends
of every link, and induces a function $\beta\colon Z\to\R$ that is injective
and increasing on every piece. Any linear order on $Z$ in which $z<z'$
whenever $\beta(z)<\beta(z')$ makes $e$ and all $a_i$ increasing.
\end{proof}
